% Einheit 2 von 2 – Die Näherungsdeutung (bank/uneigentliche-integrale/e2.jsonl, zusammenbau v0.1)
%% TODO Zeitmarke Einheit 2: Marken-Zeile nicht lesbar
%% TODO Zweigzeile Teil 1: was hier gelernt wird (Satz in Schülersprache aus dem Einheitstitel) – Einheit 2
\einheitenkopf[e2]{Einheit 2 von 2 · Die Näherungsdeutung}
\zweigzeile{Abitur LK}
\setcounter{aufgabe}{7}

%% TODO Titel: Ich-kann-Satz fehlt in der Bank (Platzhalter: Kettenname) – Nr. 8
%% TODO Anweisung: ein Satz über der ersten Teilaufgabe fehlt in der Bank – Nr. 8
\begin{aufgabe}{Die Näherungsdeutung}
%% TODO \rechenplatz{n}: Schrittzahl des Grundfalls fehlt in der Bank (nur bei fester Schreibform, 2.3 b)
\begin{teile}
\teil Die Fläche $R$ rechts von $x = 5$ ist schraffiert. Fällt sie gegenüber der Fläche links davon ins Gewicht? \\ \kreuz{fällt ins Gewicht} \\ \kreuz{vernachlässigbar klein}

\begin{ksys}[xmin=0,xmax=10,ymin=0,ymax=3]\funktion{3*exp(-\x)}{f}\flaeche{3*exp(-\x)}{5}{10}{R}\end{ksys}
\teil $F$ ist eine Stammfunktion von $f(x) = 2e^{-0{,}5x}$. Was bedeutet $F(b) - F(0)$ für $b > 0$ geometrisch?
\teil Der Graph von $f(x) = 4e^{-x}$ ist gezeichnet. Markiere die Restfläche zwischen $x = 6$ und einem großen $w$ und begründe, warum sie klein ist.

\begin{ksys}[xmin=0,xmax=10,ymin=0,ymax=4]\funktion{4*exp(-\x)}{f}\end{ksys}
\teil $f(x) = 6x \cdot e^{-0{,}4x}$ für $x \ge 0$; der Graph läuft für wachsendes $x$ gegen null. Für jede Stammfunktion $F$ von $f$ und jedes $b > 40$ gilt $F(b) - F(0) \approx ∫_0^{40} f(x)\,dx$. Deute diese Aussage geometrisch \hfill (Abitur 2022 LK).
\end{teile}
\end{aufgabe}

%% TODO Titel: Ich-kann-Satz fehlt in der Bank (Platzhalter: Kettenname) – Nr. 9
%% TODO Anweisung: ein Satz über der ersten Teilaufgabe fehlt in der Bank – Nr. 9
\begin{aufgabe}{Die Näherungsdeutung – Fehler finden, Begründen}
\begin{teile}
\teil Ich finde den Fehler: Für jede Stammfunktion $F$ von $f > 0$ und jedes $b > 20$ gilt $F(b) - F(0) \approx ∫_0^{20} f(x)\,dx$. Tom: „Also ist $F(b)$ für große $b$ ungefähr gleich $F(0)$.“
\teil Begründe: Für $f \ge 0$ ist $F(b) - F(a)$ mit $a < b$ ein Flächeninhalt.
\end{teile}
\end{aufgabe}
